\documentclass[12pt, twoside]{article}
\input{../preamble}

\fancyhead[LE]{\thepage}
\fancyhead[RO]{Markscheme (b)}
\fancyhead[LO]{La Scuola d'Italia / Dr.\ Huson / IB Math (b) \\* 17 April 2026}

\begin{document}
\subsubsection*{6.8 Quiz: Polynomials Challenge --- Markscheme (b)}

\begin{enumerate}[itemsep=0.3cm]

%--- Problem 1: f(x) = -2(x+3)(x-5) [7 marks] ---
\item[1.] [Maximum mark: 7]

Consider the function $f(x) = -2(x + 3)(x - 5)$, for $x \in \mathbb{R}$.

\medskip

\textbf{(a)(i)} Write down the $x$-intercepts. \hfill [2]

$x = -3$ and $x = 5$ \hfill \textbf{(A1)(A1)}

\medskip

\textbf{(a)(ii)} Find the coordinates of the vertex. \hfill [3]

Axis of symmetry: $x = \dfrac{-3 + 5}{2} = 1$ \hfill \textbf{(M1)}

$f(1) = -2(1 + 3)(1 - 5) = -2(4)(-4) = 32$ \hfill \textbf{(A1)}

Vertex: $(1, 32)$ \hfill \textbf{(A1)}

\medskip

\textbf{(b)} $f(x) = -2(x - h)^2 + k$. Write down $h$ and $k$. \hfill [2]

$h = 1$ \hfill \textbf{(A1)}

$k = 32$ \hfill \textbf{(A1)}

\emph{Verification: $-2(x - 1)^2 + 32 = -2(x^2 - 2x + 1) + 32
= -2x^2 + 4x - 2 + 32 = -2x^2 + 4x + 30$}

\emph{And: $-2(x + 3)(x - 5) = -2(x^2 - 2x - 15) = -2x^2 + 4x + 30$ \checkmark}

\bigskip
\hrule
\bigskip

%--- Problem 2: y = mx tangent to y = x^2 + 16 [7 marks] ---
\item[2.] [Maximum mark: 7]

The line $y = mx$ is tangent to the parabola $y = x^2 + 16$.

\medskip

\textbf{(a)} Show that at any intersection, $x^2 - mx + 16 = 0$. \hfill [1]

$x^2 + 16 = mx \implies x^2 - mx + 16 = 0$ \hfill \textbf{(A1)(AG)}

\medskip

\textbf{(b)} Find the discriminant in terms of $m$. \hfill [1]

$\Delta = (-m)^2 - 4(1)(16) = m^2 - 64$ \hfill \textbf{(A1)}

\medskip

\textbf{(c)} Find the values of $m$ for tangency. \hfill [2]

Tangent: $\Delta = 0 \implies m^2 - 64 = 0$ \hfill \textbf{(M1)}

$m = \pm 8$ \hfill \textbf{(A1)}

\medskip

\textbf{(d)} Find the coordinates of each point of tangency. \hfill [3]

For $m = 8$: $x^2 - 8x + 16 = 0 \implies (x - 4)^2 = 0
\implies x = 4$ \hfill \textbf{(M1)}

$y = 8(4) = 32$. \quad Point of tangency: $(4, 32)$ \hfill \textbf{(A1)}

For $m = -8$: $x^2 + 8x + 16 = 0 \implies (x + 4)^2 = 0
\implies x = -4$

$y = -8(-4) = 32$. \quad Point of tangency: $(-4, 32)$ \hfill \textbf{(A1)}

\newpage

%--- Problem 3: f(x) = 3x^2 + 6x - 2 [7 marks] ---
\item[3.] [Maximum mark: 7]

Consider the function $f(x) = 3x^2 + 6x - 2$.

\medskip

\textbf{(a)} Write $f(x)$ in the form $a(x - h)^2 + k$. \hfill [3]

$f(x) = 3(x^2 + 2x) - 2$ \hfill \textbf{(M1)}

$= 3(x^2 + 2x + 1 - 1) - 2 = 3(x + 1)^2 - 3 - 2$ \hfill \textbf{(M1)}

$= 3(x + 1)^2 - 5$ \hfill \textbf{(A1)}

\medskip

\textbf{(b)} Write down the coordinates of the vertex. \hfill [1]

Vertex: $(-1, -5)$ \hfill \textbf{(A1)}

\medskip

\textbf{(c)} Find the discriminant of $f(x) = 0$. \hfill [1]

$\Delta = (6)^2 - 4(3)(-2) = 36 + 24 = 60$ \hfill \textbf{(A1)}

\medskip

\textbf{(d)} State the number of $x$-intercepts and justify. \hfill [2]

Since $\Delta = 60 > 0$, $f(x) = 0$ has two distinct real solutions. \hfill \textbf{(A1)}

Therefore the graph has two $x$-intercepts. \hfill \textbf{(R1)}

\bigskip
\hrule
\bigskip

%--- Problem 4: p(x) = 2x^3 - 3x^2 - 11x + 6 [8 marks] ---
\item[4.] [Maximum mark: 8]

Let $p(x) = 2x^3 - 3x^2 - 11x + 6$.

\medskip

\textbf{(a)} Show that $(x - 3)$ is a factor of $p(x)$. \hfill [2]

$p(3) = 2(27) - 3(9) - 11(3) + 6$ \hfill \textbf{(M1)}

$= 54 - 27 - 33 + 6 = 0$

Since $p(3) = 0$, $(x - 3)$ is a factor by the
factor theorem. \hfill \textbf{(A1)(AG)}

\medskip

\textbf{(b)} Express $p(x) = (x - 3)(ax^2 + bx + c)$. \hfill [2]

By polynomial division: \hfill \textbf{(M1)}

$2x^3 - 3x^2 - 11x + 6 = (x - 3)(2x^2 + 3x - 2)$ \hfill \textbf{(A1)}

\emph{i.e.\ $a = 2$, $b = 3$, $c = -2$.}

\medskip

\textbf{(c)} Factorise $p(x)$ completely. \hfill [2]

$2x^2 + 3x - 2 = (2x - 1)(x + 2)$ \hfill \textbf{(M1)}

$p(x) = (x - 3)(2x - 1)(x + 2)$ \hfill \textbf{(A1)}

\medskip

\textbf{(d)} Write down all solutions to $p(x) = 0$. \hfill [1]

$x = 3, \quad x = \dfrac{1}{2}, \quad x = -2$ \hfill \textbf{(A1)}

\medskip

\textbf{(e)} Write down the $y$-intercept. \hfill [1]

$p(0) = 6$, so the $y$-intercept is $(0, 6)$. \hfill \textbf{(A1)}

\end{enumerate}
\end{document}
